

\section{Related Work}\label{sec:related}


\noindent \textbf{Multi-Agent Systems.} Multi-Agent Systems that leverage multiple LLM-based agents communicating and collaborating, are increasingly employed to tackle complex, multi-step challenges \cite{tran2025multi}. Existing research and industrial applications have demonstrated the effectiveness and robustness of these systems across diverse domains, such as code generation \cite{barbarroxa2024benchmarking}, reasoning \cite{debate1-mit}, question answering \cite{das2023enabling}, and decision making \cite{sun2024llm}. Two critical factors in designing MAS are agent composition and the collaboration structure \cite{tran2025multi}. We categorize these factors as relating to the functional and structural properties of MAS.


% \noindent \textbf{Agent Construction in MAS.} As the core components in MAS, the design and construction of agents directly determine the system's overall functional capabilities. Most existing works on MAS employs either manually designed prompts or prompts generated by LLMs to construct agents tailored to specific application scenarios. For example, \citet{das2023enabling} constructed agent teams with specific roles (e.g., CEO, programmer, tester) using hand-crafted instruction prompts for software development. \citet{multi-persona} utilized LLMs to generate role-specific prompts for agents in response to task-specific queries. However, both manual prompt design and LLM-driven agent generation rely heavily on human prior knowledge, requiring empirical verification through trial-and-error. Although previous studies have explored agent optimization through fine-tuning \cite{metatool, api_bank} or prompt-tuning techniques \cite{dsp, g_retriever}, these methods do not explicitly address the optimization of agents within MAS involving multi-step collaborative processes.


\noindent \textbf{Agent Construction in MAS.} 
% As the core components in MAS, the design and construction of agents directly determine the system's overall functional capabilities. 
The construction of agents directly determines the system's overall functional capabilities. 
Most existing works on MAS employ either manually designed prompts or prompts generated by LLMs to construct agents tailored to specific application scenarios. For example, \citet{das2023enabling} constructed agent teams with specific roles using hand-crafted instruction prompts for software development. \citet{multi-persona} utilized LLMs to generate role-specific prompts for agents in response to task-specific queries. However, both manual prompt design and LLM-driven agent generation rely heavily on human prior knowledge, requiring empirical verification through trial-and-error. Although previous studies have explored agent optimization through fine-tuning \cite{api_bank} or prompt-tuning techniques \cite{dsp}, these methods do not explicitly address the optimization of agents within MAS involving multi-step collaborative processes.


\noindent \textbf{Collaboration Structure in MAS.} Another crucial aspect of MAS is the collaboration structure, which defines how candidate agents collaborate and communicate to ultimately resolve the given task. Existing research has proposed various structures, such as centralized \cite{jiang2023llm}, decentralized \cite{liang2023encouraging}, and hierarchical \cite{chatllm-network} approaches. However, these structures are typically manually designed for specific task categories and remain static throughout the collaboration process. A recent work, DyLAN \cite{dylan}, represents an initial effort in dynamically optimizing the collaboration structure within MAS. However, its focus is limited to
optimizing agent team composition, relying on an unsupervised metric called ``Agent Importance Score'' computed using heuristic rules combined with LLM-based judgments. In contrast, OMAC comprehensively examines both functional and structural optimization of MAS in a joint supervised manner, addressing more fine-grained optimization of collaboration structures.

Due to the space limitation, we've included expanded discussion of related works on MAS optimization in Appendix~\ref{sec:exp_related}.





% Another key aspect of MAS is the collaboration structure, which defines how agents collaborate and communicate to ultimately resolve the given task. Existing research has proposed various structures, including centralized \cite{}, decentralized \cite{}, and hierarchical \cite{@InProceedingsChen_2022_CVPR} approaches. However, these structures are typically manually defined for specific task types and remain static throughout the collaboration process. A recent work, DyLAN \cite{}, represents an initial effort in dynamically optimizing the collaboration structure within MAS. However, its focus is limited to optimizing agent team composition (i.e., selecting participating agents), relying on an unsupervised metric ("Agent Importance Score") derived from heuristic rules and LLM-based judgments. In contrast, our work comprehensively addresses both functional and structural optimization in MAS using a supervised approach, including finer-grained optimization of the collaboration structure itself.



% As the core components of MAS, the design and construction of agents directly determine the functional scope of these systems. Most existing research on multi-agent collaboration employs either manually designed prompts or prompts generated by LLMs to construct agents tailored to specific application scenarios. For example, \citet{das2023enabling} constructed agent teams with predefined roles (e.g., CEO, programmer, tester) using manually crafted instruction prompts for software development tasks. Similarly, \citet{multi-persona} utilized LLMs to generate role-specific prompts for agents in response to task-specific queries. However, both manual prompt design and LLM-driven prompt generation heavily rely on human prior knowledge, requiring empirical verification through trial-and-error. Although previous studies have explored agent optimization through fine-tuning \cite{metatool, gpt4_tools, api_bank} or prompt-tuning techniques \cite{retroformer, prompt_agent, dsp, g_retriever}, these methods do not explicitly address the optimization of agents within MAS involving multi-step collaborative processes.


\endinput

We briefly introduce the related works in conversational recommendation, reinforcement learning, and graph learning.

% The success of a recommendation system hinges on offering the relevant items of user interest accurately and timely. At beginning, recommendation systems are largely built on the collaborative filtering hypothesis to infer a distributed representation of the user profile. Representative models include matrix factorization~\cite{fastMF} and factorization machines~\cite{FM,NFM}. However, by nature, these approaches suffer from two intrinsic problems. The first one is the inability of capturing user dynamic preferences with the strict assumption that a user's interest is static over the long-term horizon~\cite{song2019session}. The second problem is the weak explainability as the user preference representation is only a continuous vector. Later works try to introduce Markov models~\cite{rendle2010factorizing} and multi-arm bandit methods~\cite{wu2018learning} to solve the dynamic problem but the explainability still remains to be unsatisfactory.
% \par
%Traditional recommendation systems lack transparency, leaving users unable to understand why certain items are recommended. Additionally, fully automated methods are often affected by the filter bubble effect, limiting the recommendations to items similar to those that have been positively evaluated in the past. Interactive control approaches can enhance the transparency and acceptability of the system. \textbf{Interactive recommendation} is a subfield of recommendation research that focuses on incorporating user interaction into the recommendation process in some way, either through explicit feedback (e.g. rating items), implicit feedback (e.g. clicking on items), or contextual information (e.g. time, location). % Choice-based recommendation is one type of interactive recommendation system. They mainly adopt reinforcement learning techniques with the aim of creating dynamic and personalized experiences for the user, adapting recommendations to their evolving preferences and contextual factors through ongoing adjustment."

\subsection{Conversational recommendation system} (CRSs) is a novel solution to recommendation that leverage natural language to effectively elicit dynamic user preferences that align with their real needs through multiple rounds of real-time interaction. CRS is considered to be a cutting-edge discipline that incorporates dialogue systems, recommendation systems, and interactive systems \cite{gao2021advances}.
% According to different functions, a CRS can be decoupled into five modules: (1) Question-based user preference elicitation. (2) Multi-turn conversational recommendation strategies. (3) Dialogue understanding and generation. (4) Exploitation-exploration trade-offs. (5) Evaluation and user simulation. 
According to the focus on different functions and settings, existing CSR methods can be roughly divided into two types: dialogue-based recommendation \cite{nips18/DeepConv,zhou2020topicguided,chen-etal-2019-towards,10.1145/3488560.3498514,wu2019proactive} and multi-round conversational recommendation (MCR) \cite{lei2020interactive,deng2021unified,xu2021adapting,10.1145/3523227.3546755,gao2022kuairec,li2020seamlessly}. In this work, we focus on the MCR setting.
% \textbf{Question-based recommendation methods}. 

MCR is considered to be the most realistic setting in CRS. Unlike dialogue-based recommenders that need to extract information or generate responses through raw natural language \cite{10.1145/3534678.3539382}, MCR focuses on the core logic of the interaction strategy which involves asking questions \cite{10.1145/3397271.3401180,10.1145/3397271.3401180,10.1145/3477495.3532077} and making recommendations.
% In MCR, the CRS is typically formulated as a multi-step decision-making process involving two key tasks: asking questions and making recommendations. The aim of asking questions is to elicit user preferences on certain attributes and then narrow down the potiential recommendation candidates. 
The traditional MCR setting allows users to select only one preferred attribute value at a time, which restricts users' expression in the interaction. 
To overcome this issue, \citet{zhang2022multiple} propose the MIMCR setting, where a user is allowed to select multiple options for a certain attribute. Though effective, they follow the recommendation philosophy in MCR to directly filter out the items that the user has not mentioned by attributes, which leads to failure as users may not be sure what they want precisely. In our proposed VPMCR setting, we specifically consider users' vague preferences and adjust the recommendation mechanism to consider the items with unmentioned attributes, which better reflect users' needs.

\subsection{RL-based Recommendation} Reinforcement Learning (RL) is a type of Machine Learning. It considers how an agent (e.g., a machine) should automatically make decisions within a specific context to pursue a long-term goal. The agent learns and adjusts its policy based on the reward feedback (i.e., reinforcement signals) given by the environment. Recently, RL has shown its effectiveness in recommendation \cite{afsar2022reinforcement,deffayet2023offline,gao2023alleviating}. As fitting user interest is not a bottleneck for now, recommenders care more about users' long-term satisfaction \cite{ResAct,10.1145/3534678.3539040,wang2022best}. For instance, \citet{10.1145/3488560.3498526} use RL to generate the proper questions that can maximally make the system help users search desired products. \citet{gao2022cirs} integrate causal inference into offline RL to maximize users' long-term satisfaction by removing filter bubbles. \citet{10.1145/3534678.3539095} propose an RL-based dispatching solution for ride-hailing platforms that can conduct robust and efficient on-policy learning and inference while being adaptable for full-scale deployment. In this work, we use RL to learn a policy that can automate question-asking and item recommendation.


\subsection{Graph-based Recommendation} Graph-based recommender systems have drawn a lot of research attention \cite{10.1145/3534678.3539452,10.1145/3534678.3539458,10.1145/3447548.3467384,10.1145/3447548.3467384}. By arranging the various entities (e.g., users, items, and attributes) in a heterogeneous graph, we can leverage lots of properties in modeling the collaborative signals. In CRS, the knowledge graph is utilized in enriching the system with additional knowledge \cite{lei2020interactive,xu2020user,zhou2020improving,moon-2019-opendialkg,zhou2020interactive}. For example, to better understand concepts that a user mentioned, \citet{zhou2020improving} propose to incorporate two external knowledge graphs (KGs): a word-oriented KG providing relations (e.g., synonyms, antonyms, or co-occurrence) between words and an item-oriented KG carrying structured facts regarding the attributes of items. With the increasing of nodes, the computational overhead is too large to satisfy the requirement of real-time interaction. Hence, we propose a pruning strategy to overcome this work.